\paragraph{Detalle de videojuego}
Desde la vista de detalle de un videojuego, el usuario puede añadir o quitar el título de dos listas distintas: la lista de videojuegos jugados (playedList) y
la lista de deseos (wishList). Para añadir el videojuego a la lista, se realiza la petición POST al endpoint 'wishList/' o 'playedList/'(RF-3.3,
RF-3.1). Si el videojuego ya forma parte de una de estas listas, 
se envía una petición DELETE a wishlist/{id} o playedList/{id}, siendo {id} el identificador del videojuego(RF-3.2, RF-3.4).
Cuando se accede a la vista de detalle, el sistema realiza automáticamente una consulta para comprobar si 
el videojuego ya está incluido en alguna de las listas. En función de esa comprobación, se modifica el icono del botón (coloreándolo de negro si ya está añadido) 
y se muestra un tooltip que indica la acción que se realizará al pulsarlo.

En el backend, se devuelve únicamente un indicador de éxito o fallo según el resultado de la operación de añadido o eliminación.

El flujo de interración se observa en la figura \ref{FIG:DETALLE}

\paragraph{Sección 'Mis listas'}
Desde la pantalla de listas personalizadas, se ofrecen dos vistas, una para crear la lista y otra para visualizar las listas. A la hora de 
visualizar las listas, el componente CustomListLayout realiza una solictud al endpoint 'myLists/'. Una vez que se obtiene la respuesta, este componente, 
llama a otro componente, encangado de la visualización de las listas. Para cada lista, se pueden realizar dos acciones 'Eliminar' o 'Editar' (RF-3.5, RF-3.6). Para 
aplicar esas acciones, se realiza la petición PUT o DELETE al endpoint 'api/v1/myLists/<str:pk>/', donde pk es el identificador de la lista. 

A la hora de crear listas, el usuario puede definir un nombre y una descripción para cada lista personalizada. 
El formulario de creación de listas realiza una solicitud POST al endpoint 'api/v1/myLists/' enviando los datos 
proporcionados por el usuario. Antes de enviar los datos, se comprueba desde el front que se ha proporcionado un nombre y 
videojuegos asociados a esa lista(RF-3.5.1). Por último, en el backend, se valida la información y, si es correcta, se almacena la nueva 
lista en la base de datos.

El flujo de interración se observa en la figura \ref{FIG:LAS}

\begin{figure}[Componentes implicados para la sección 'Mis listas']{FIG:LAS}{ En este diagrama se puede ver el proceso completo de la gestión de listas}
    \image{13cm}{}{listaLayout}
\end{figure}

\paragraph{Sección 'Played' y 'Wishlist'}
Las pantallas de 'Played' y 'Wishlist' se implementan de manera similar realizando un petición GET al endpoint '/wishlist' o '/playedList'. La petición 
la realiza el componente WishlistLayout y para la visualización de los videojuegos se reutiliza el componente VideogameList.

Del lado del backend, en cada view set se obtienen los videojuegos asociados en función del atributo 'IsPlayedList' de la tabla 'Wishlist'.

El flujo de interración se observa en la figura \ref{FIG:WISH}

\begin{figure}[Componentes implicados para la sección 'Amistades']{FIG:WISH}{ En este diagrama se puede ver el proceso completo de la gestión de amistades}
    \image{13cm}{}{wishlist}
\end{figure}